\section{Aspects of the Data}

\subsection{Selection of the Measure \label{measure}}

Selecting an appropriate measure was an important step in assembling the data. The OECD's \textit{National Accounts} offers 27 measures, of which 20 are potentially relevant for the levels of the variables studied here. The remaining seven concern growth or population.

BKK do not specify their chosen measure. They state only that the variables are measured in real terms, except for the net exports-to-output ratio, which uses current prices.

We used three steps to select a comparable measure:

\begin{enumerate}
  \item We downloaded output data from the OECD website in every available measure.
  \item We used R to calculate output correlations between the United States and the ten other BKK countries for each measure.
  \item We compared the results with BKK's and selected the closest comparable measure.
\end{enumerate}

Over the full period from 1947:1 to 2015:2, nine measures provide enough observations to calculate output correlations between the United States and the other sample countries. Restricting the comparison to BKK's period, 1970:1 to 1990:2, leaves four measures. \Cref{tab:measure} reports the resulting US output correlations alongside those from BKK and ACZ.

\begin{table}[htbp]
  \centering
\begin{adjustbox}{max width=\textwidth}
\begin{tabular}{lrrrrrrrrrr}
  \toprule
 & AUS & AUT & CAN & CHE & DEU & FRA & GBR & ITA & JPN & USA \\
  \midrule
  \cite{Backus:1993aa} & .51 & .38 & .76 & .42 & .69 & .41 & .55 & .41 & .60 & 1.00 \\
\cite{Ambler:2004aa} & .40 & .32 & .77 & .35 & .50 & .27 & .68 & .33 & .39 & 1.00 \\
CARSA & .00 & .08 & .29 & .11 & .29 & -.14 & -.01 & -.04 & .22 & 1.00 \\
  CPCARSA & .27 & -.00 & .56 & .26 & .41 & .11 & .30 & .24 & .38 & 1.00 \\
  VOBARSA & .35 & .29 & .71 & .41 & .66 & .39 & .62 & .37 & .60 & 1.00 \\
  VPVOBARSA & .35 & .29 & .71 & .41 & .66 & .39 & .62 & .37 & .60 & 1.00 \\
   \bottomrule
\end{tabular}
\end{adjustbox}
\caption{US Output Correlations under Four Measures, 1970:1--1990:2}
\label{tab:measure}
\end{table}

We exclude CARSA and CPCARSA because they measure output at current prices, whereas BKK use real values. Their correlations also differ substantially from BKK's.

VOBARSA and VPVOBARSA produce identical correlations and are the closest to both BKK's and ACZ's estimates. They differ only in scale: VOBARSA is expressed in national currency, while VPVOBARSA is expressed in US dollars. Scaling a series by a positive constant leaves its correlations unchanged. For our analysis of HP-filtered log values, a fixed currency conversion adds only a constant to the log series and therefore leaves the cyclical component unchanged. We choose VPVOBARSA so that output levels are also expressed in a common unit.

\subsection{Civilian Employment Series \label{app:employment}}

BKK obtain quarterly civilian employment data from the OECD's \textit{Main Economic Indicators}. At the time of data collection, these series were unavailable on the OECD website, which instead provided quarterly employed-population data through its \textit{Short-Term Labour Market Statistics}. As \cref{tab:emptimespan} shows, some of these series are too short for a useful comparison with BKK. The Federal Reserve Bank of St. Louis website\footnote{\url{https://research.stlouisfed.org/}} provides the \textit{Main Economic Indicators} series for France, the United Kingdom and Italy, allowing us to extend their coverage. We could not find comparable series for Switzerland or the European aggregate, so we exclude them from the employment comovement calculations.

\begin{table}[ht]
\centering
\begin{adjustbox}{max width=\textwidth}
\begin{tabular}{rccccccccccc}
  \toprule
Country & AUS & AUT & CAN & FRA & DEU & ITA & JPN & CHE & GBR & USA & EA19 \\
  \midrule
  Last  & 2015:2 & 2015:1 & 2015:2 & 2015:1 & 2015:1 & 2015:1 & 2015:2 & 2015:1 & 2015:1 & 2015:2 & 2015:1 \\
  First  & 1967:1 & 1969:1 & 1955:1 & 2003:1 & 1962:1 & 1998:1 & 1955:1 & 1998:2 & 1999:2 & 1955:1 & 2005:1 \\
   \bottomrule
\end{tabular}
\end{adjustbox}
\caption{Coverage of Employed Population Series from \textit{Short-Term Labour Market Statistics}}
\label{tab:emptimespan}
\end{table}

\Cref{fig:emp} compares the employment series from the two sources. Their similar behavior supports splicing them to obtain longer series for France, the United Kingdom and Italy. We rescale the \textit{Short-Term Labour Market Statistics} series so that it matches the \textit{Main Economic Indicators} series at the splice point, then use it to extend the earlier series. \Cref{tab:timespan} reports the coverage of the resulting series.

\begin{figure}[htb]
\centering
\begin{subfigure}{0.46\textwidth}
  \centering
\includegraphics[width=\linewidth]{\generatedfigures/fraemployment.png}
  \caption{France}
  \label{fig:fraemp}
\end{subfigure}
\qquad
\begin{subfigure}{0.46\textwidth}
  \centering
\includegraphics[width=\linewidth]{\generatedfigures/gbremployment.png}
  \caption{United Kingdom}
  \label{fig:gbremp}
\end{subfigure}

\medskip

\begin{subfigure}{0.62\textwidth}
  \centering
\includegraphics[width=\linewidth]{\generatedfigures/itaemployment.png}
  \caption{Italy}
  \label{fig:itaemp}
\end{subfigure}
\caption{Civilian Employment Series}
\label{fig:emp}
\end{figure}
